\chapter{Giỏi Không Đồng Nghĩa Với Kiếm Được Việc}
\markboth{Giỏi Không Đồng Nghĩa Với Kiếm Được Việc}{Being Good Does Not Automatically Get You a Job}

\section{Giỏi Trên Giấy Chưa Chắc Giỏi Trong Phòng Phỏng Vấn}
\begin{truyen}{Giỏi Trên Giấy Chưa Chắc Giỏi Trong Phòng Phỏng Vấn}{Good on Paper Does Not Mean Good in the Interview Room}
\chuhoa{C}{ó người bảng điểm rất đẹp nhưng vào phỏng vấn nói lan man, thiếu rõ, và không cho thấy mình hiểu công việc.}
\textit{(Some people have excellent transcripts but speak vaguely in interviews and do not show that they understand the work.)}

Nhà tuyển dụng không chỉ nhìn thành tích.
\textit{(Employers do not only look at achievements.)}

Họ nhìn xem em có thể biến kiến thức thành hành động hay không.
\textit{(They look at whether you can turn knowledge into action.)}

Người trẻ hay tưởng học giỏi là đủ.
\textit{(Young people often think being academically strong is enough.)}

Nhưng đời cần thêm cách trình bày, cách ứng xử, và cách chứng minh.
\textit{(But life also needs communication, conduct, and proof.)}

Một năng lực không biết cách xuất hiện thì rất dễ bị bỏ qua.
\textit{(An ability that does not know how to show up is easy to overlook.)}

\end{truyen}

\begin{baihoc}
Học giỏi rất quý. Nhưng muốn có việc, em còn phải học cách nói rõ mình làm được gì và vì sao mình hợp.
\textit{(Studying well is valuable. But to get work, you must also learn how to explain what you can do and why you fit.)}
\end{baihoc}

\begin{ghinhoanh}
\textbf{Short English to remember:}

Studying well is valuable.

But to get work, you must also learn how to explain what you can do and why you fit.

\end{ghinhoanh}

\begin{tuhocnhanh}
\textbf{Cau hoi tu hoc / Self-study questions}

1. Van de chinh la gi? \textit{What is the main problem?}

2. Bai hoc thuc te la gi? \textit{What is the practical lesson?}

3. Mot cau tieng Anh em muon nho la cau nao? \textit{Which English sentence do you want to remember?}
\end{tuhocnhanh}
\ngancach

\section{Kiến Thức Không Ứng Dụng Được Sẽ Bị Nghi Ngờ}
\begin{truyen}{Kiến Thức Không Ứng Dụng Được Sẽ Bị Nghi Ngờ}{Knowledge Without Use Will Be Questioned}
\chuhoa{N}{hiều bạn trẻ học được rất nhiều nhưng không làm ra sản phẩm nào, không có ví dụ nào, không có việc gì chỉ rõ mình từng làm được.}
\textit{(Many young people learn a lot but produce nothing, show no examples, and cannot point to work they have actually done.)}

Khi đó nhà tuyển dụng phải tin vào lời em nói nhiều hơn tin vào dấu vết công việc.
\textit{(Then an employer must trust your words more than any evidence of work.)}

Điều đó làm em yếu đi trong cạnh tranh.
\textit{(That weakens you in competition.)}

Kiến thức mạnh hơn khi nó để lại dấu vết: bài viết, dự án, sản phẩm, kỹ năng dùng được.
\textit{(Knowledge becomes stronger when it leaves traces: writing, projects, products, and usable skill.)}

\end{truyen}

\begin{baihoc}
Muốn người khác tin mình có năng lực, hãy tạo ra bằng chứng nhỏ nhưng thật. Đừng chỉ nói em biết.
\textit{(If you want others to believe in your ability, create small but real evidence. Do not only say that you know.)}
\end{baihoc}

\begin{ghinhoanh}
\textbf{Short English to remember:}

If you want others to believe in your ability, create small but real evidence.

Do not only say that you know.

\end{ghinhoanh}

\begin{tuhocnhanh}
\textbf{Cau hoi tu hoc / Self-study questions}

1. Van de chinh la gi? \textit{What is the main problem?}

2. Bai hoc thuc te la gi? \textit{What is the practical lesson?}

3. Mot cau tieng Anh em muon nho la cau nao? \textit{Which English sentence do you want to remember?}
\end{tuhocnhanh}
\ngancach

\section{Biết Nhiều Không Bằng Làm Xong}
\begin{truyen}{Biết Nhiều Không Bằng Làm Xong}{Knowing More Is Not Better Than Finishing}
\chuhoa{C}{ó người học rất rộng nhưng việc nào cũng để dở.}
\textit{(Some people study widely but leave every task unfinished.)}

Ngoài đời, người hoàn thành ổn định thường được tin hơn người biết rất nhiều nhưng hay bỏ lửng.
\textit{(In real life, people who finish reliably are often trusted more than people who know a lot but leave things hanging.)}

Hoàn thành là một đức tính ít hào nhoáng nhưng rất mạnh.
\textit{(Finishing is a quiet virtue, but a powerful one.)}

Người trẻ muốn có việc phải học cách đi tới cuối, không chỉ hứng khởi ở đầu.
\textit{(Young people who want work must learn to go to the end, not only feel excited at the start.)}

\end{truyen}

\begin{baihoc}
Thị trường thường thưởng cho người làm xong và làm được, không chỉ cho người nói hay hay biết nhiều.
\textit{(The market often rewards people who finish and deliver, not only those who speak well or know a lot.)}
\end{baihoc}

\begin{ghinhoanh}
\textbf{Short English to remember:}

The market often rewards people who finish and deliver, not only those who speak well or know a lot.

\end{ghinhoanh}

\begin{tuhocnhanh}
\textbf{Cau hoi tu hoc / Self-study questions}

1. Van de chinh la gi? \textit{What is the main problem?}

2. Bai hoc thuc te la gi? \textit{What is the practical lesson?}

3. Mot cau tieng Anh em muon nho la cau nao? \textit{Which English sentence do you want to remember?}
\end{tuhocnhanh}
\ngancach

\section{Giỏi Nhưng Khó Làm Việc Cùng Là Một Điểm Trừ Lớn}
\begin{truyen}{Giỏi Nhưng Khó Làm Việc Cùng Là Một Điểm Trừ Lớn}{Skilled but Hard to Work With Is a Major Problem}
\chuhoa{C}{ó người rất thông minh nhưng phản ứng gay gắt khi bị góp ý, coi thường người khác, và làm không khí chung nặng nề.}
\textit{(Some people are very smart but react harshly to feedback, look down on others, and make the team atmosphere heavy.)}

Ban đầu họ vẫn được giữ vì năng lực.
\textit{(At first they may still be kept because of ability.)}

Nhưng về lâu dài, nhiều nơi không muốn trả giá quá lớn cho một người giỏi mà gây hao tổn tập thể.
\textit{(But in the long run, many workplaces do not want to pay too high a price for one skilled person who drains the whole group.)}

Cách làm việc cùng người khác là một phần của năng lực, không phải chuyện phụ.
\textit{(The way you work with others is part of competence, not a side issue.)}

\end{truyen}

\begin{baihoc}
Muốn có chỗ đứng bền, người trẻ phải rèn cả tài và nết làm việc. Chỉ có một vế thì rất dễ chông chênh.
\textit{(To keep a stable place, young people must develop both skill and work character. Having only one side is shaky.)}
\end{baihoc}

\begin{ghinhoanh}
\textbf{Short English to remember:}

To keep a stable place, young people must develop both skill and work character.

Having only one side is shaky.

\end{ghinhoanh}

\begin{tuhocnhanh}
\textbf{Cau hoi tu hoc / Self-study questions}

1. Van de chinh la gi? \textit{What is the main problem?}

2. Bai hoc thuc te la gi? \textit{What is the practical lesson?}

3. Mot cau tieng Anh em muon nho la cau nao? \textit{Which English sentence do you want to remember?}
\end{tuhocnhanh}
\ngancach

\section{Đi Xin Việc Cũng Là Một Kỹ Năng}
\begin{truyen}{Đi Xin Việc Cũng Là Một Kỹ Năng}{Job Seeking Is Also a Skill}
\chuhoa{N}{hiều người nghĩ chỉ cần giỏi là người ta sẽ tự tìm tới.}
\textit{(Many people think that if they are good enough, others will automatically find them.)}

Điều đó hiếm khi đúng lúc mới bắt đầu.
\textit{(That is rarely true at the beginning.)}

Viết CV, gửi hồ sơ, viết thư ngắn gọn, trả lời email, phỏng vấn, và theo dõi sau buổi gặp đều là kỹ năng phải học.
\textit{(Writing a CV, sending applications, writing short emails, interviewing, and following up are all skills that must be learned.)}

Người trẻ càng sớm hiểu đi xin việc cũng là một nghề phụ tạm thời, họ càng bớt bị động.
\textit{(The sooner young people understand that job seeking is a temporary craft of its own, the less passive they become.)}

\end{truyen}

\begin{baihoc}
Đừng chỉ học để làm việc. Hãy học cả cách bước vào một công việc. Cánh cửa cũng có kỹ năng riêng của nó.
\textit{(Do not only learn how to do the job. Learn how to enter the job. The door has its own skill too.)}
\end{baihoc}

\begin{ghinhoanh}
\textbf{Short English to remember:}

Do not only learn how to do the job.

Learn how to enter the job.

The door has its own skill too.

\end{ghinhoanh}

\begin{tuhocnhanh}
\textbf{Cau hoi tu hoc / Self-study questions}

1. Van de chinh la gi? \textit{What is the main problem?}

2. Bai hoc thuc te la gi? \textit{What is the practical lesson?}

3. Mot cau tieng Anh em muon nho la cau nao? \textit{Which English sentence do you want to remember?}
\end{tuhocnhanh}
\ngancach

\section{Hồ Sơ Có Dấu Vết}
\begin{truyen}{Hồ Sơ Có Dấu Vết}{A Portfolio Leaves Traces}
\chuhoa{C}{ó người không nói mình giỏi nhiều, nhưng mở máy ra là có bài viết, sản phẩm, và dự án để cho xem.}
\textit{(Some people do not talk much about being capable, but when they open their laptop there are articles, products, and projects to show.)}

Nhiều bạn trẻ tưởng nói hay về năng lực là đủ.
\textit{(Many young people think speaking well about ability is enough.)}

Nhưng bằng chứng nhỏ mà thật thường thuyết phục hơn lời giới thiệu lớn.
\textit{(But small real evidence persuades more than grand introductions.)}

Một dấu vết cụ thể có thể mở cửa nhanh hơn rất nhiều câu tự khen.
\textit{(One concrete trace can open doors faster than many self-promoting lines.)}

\end{truyen}

\begin{baihoc}
Muốn kiếm việc, đừng chỉ tích kiến thức. Hãy tích cả bằng chứng.
\textit{(If you want work, do not collect only knowledge. Collect evidence too.)}
\end{baihoc}

\begin{ghinhoanh}
\textbf{Short English to remember:}

If you want work, do not collect only knowledge.

Collect evidence too.

\end{ghinhoanh}

\begin{tuhocnhanh}
\textbf{Cau hoi tu hoc / Self-study questions}

1. Van de chinh la gi? \textit{What is the main problem?}

2. Bai hoc thuc te la gi? \textit{What is the practical lesson?}

3. Mot cau tieng Anh em muon nho la cau nao? \textit{Which English sentence do you want to remember?}
\end{tuhocnhanh}
\ngancach

\section{Bẫy Thực Tập Không Lối Ra}
\begin{truyen}{Bẫy Thực Tập Không Lối Ra}{The Trap of Endless Unpaid Internships}
\chuhoa{C}{ó người đi hết chỗ thực tập này sang chỗ khác mà vị trí vẫn mờ, lương vẫn không có.}
\textit{(Some move from one internship to another while the role stays unclear and pay never appears.)}

Họ tưởng cứ chịu khó ôm việc miễn phí thật lâu thì sớm muộn cũng được cất nhắc.
\textit{(They think if they carry unpaid work long enough, they will eventually be lifted up.)}

Nhưng có nơi chỉ quen dùng sức trẻ giá rẻ mà không hề có đường đi rõ cho em.
\textit{(But some places simply get used to cheap young labor and offer no clear path.)}

Người trẻ cần học nhìn xem nơi đó đang dạy mình thật hay chỉ đang dùng mình khéo.
\textit{(Young people need to ask whether a place is truly teaching them or merely using them well.)}

\end{truyen}

\begin{baihoc}
Không phải cơ hội nào cũng đáng trả bằng quá nhiều thời gian miễn phí.
\textit{(Not every opportunity is worth paying for with too much unpaid time.)}
\end{baihoc}

\begin{ghinhoanh}
\textbf{Short English to remember:}

Not every opportunity is worth paying for with too much unpaid time.

\end{ghinhoanh}

\begin{tuhocnhanh}
\textbf{Cau hoi tu hoc / Self-study questions}

1. Van de chinh la gi? \textit{What is the main problem?}

2. Bai hoc thuc te la gi? \textit{What is the practical lesson?}

3. Mot cau tieng Anh em muon nho la cau nao? \textit{Which English sentence do you want to remember?}
\end{tuhocnhanh}
\ngancach

\section{Kết Nối Không Phải Xin Xỏ}
\begin{truyen}{Kết Nối Không Phải Xin Xỏ}{Networking Is Not Begging}
\chuhoa{N}{hiều người giỏi nhưng ngại kết nối vì sợ bị nghĩ là nhờ vả hoặc đánh bóng mình.}
\textit{(Many capable people avoid connection because they fear looking needy or self-promoting.)}

Họ tưởng đi một mình mới là sạch và đáng trọng.
\textit{(They think walking alone is the purest and most respectable path.)}

Nhưng công việc thật vận hành qua con người, sự tin cậy, và những cuộc gặp đúng lúc.
\textit{(But real work moves through people, trust, and well-timed conversations.)}

Kết nối đúng không làm em thấp đi.
\textit{(Healthy networking does not lower you.)}

Nó giúp giá trị của em đi đến đúng tai đúng mắt.
\textit{(It helps your value reach the right ears and eyes.)}

\end{truyen}

\begin{baihoc}
Kiếm việc không chỉ là học giỏi. Còn là học cách bước vào mạng người thật.
\textit{(Getting work is not only about being good. It is also about entering real human networks.)}
\end{baihoc}

\begin{ghinhoanh}
\textbf{Short English to remember:}

Getting work is not only about being good.

It is also about entering real human networks.

\end{ghinhoanh}

\begin{tuhocnhanh}
\textbf{Cau hoi tu hoc / Self-study questions}

1. Van de chinh la gi? \textit{What is the main problem?}

2. Bai hoc thuc te la gi? \textit{What is the practical lesson?}

3. Mot cau tieng Anh em muon nho la cau nao? \textit{Which English sentence do you want to remember?}
\end{tuhocnhanh}
\ngancach

\section{Trượt Phỏng Vấn Đầu}
\begin{truyen}{Trượt Phỏng Vấn Đầu}{Failing the First Interview}
\chuhoa{C}{uộc phỏng vấn đầu tiên làm nhiều người trẻ vấp ngay ở những chỗ mình tưởng đơn giản.}
\textit{(The first interview makes many young people stumble on things they thought were simple.)}

Họ tưởng học giỏi rồi thì nói về mình cũng sẽ tự trôi.
\textit{(They think that if they studied well, speaking about themselves will naturally go well too.)}

Nhưng trả lời ngắn gọn, đúng trọng tâm, và cho thấy mình hiểu công việc là một kỹ năng phải luyện.
\textit{(But answering briefly, clearly, and with real understanding is a skill that must be practiced.)}

Trượt lần đầu không đáng xấu hổ.
\textit{(Failing the first time is not shameful.)}

Nó chỉ phơi ra chỗ em cần tập.
\textit{(It only reveals what you need to practice.)}

\end{truyen}

\begin{baihoc}
Phỏng vấn dở không nói em vô dụng. Nó nói em còn phải học thêm một loại ngôn ngữ của công việc.
\textit{(A poor interview does not say you are useless. It says you still need to learn one more language of work.)}
\end{baihoc}

\begin{ghinhoanh}
\textbf{Short English to remember:}

A poor interview does not say you are useless.

It says you still need to learn one more language of work.

\end{ghinhoanh}

\begin{tuhocnhanh}
\textbf{Cau hoi tu hoc / Self-study questions}

1. Van de chinh la gi? \textit{What is the main problem?}

2. Bai hoc thuc te la gi? \textit{What is the practical lesson?}

3. Mot cau tieng Anh em muon nho la cau nao? \textit{Which English sentence do you want to remember?}
\end{tuhocnhanh}
\ngancach

\section{Học Công Cụ, Không Chỉ Học Lý Thuyết}
\begin{truyen}{Học Công Cụ, Không Chỉ Học Lý Thuyết}{Learn Tools, Not Theory Alone}
\chuhoa{C}{ó người biết khái niệm rất nhiều nhưng đụng vào phần mềm, quy trình, hay công cụ thật thì lúng túng.}
\textit{(Some people know many concepts, but become awkward when touching real software, process, or tools.)}

Họ tưởng nền lý thuyết đẹp sẽ tự đổi thành khả năng làm việc.
\textit{(They think a strong theory base will automatically turn into working ability.)}

Nhưng thị trường thường nhìn người dùng được công cụ và giải được việc cụ thể trước tiên.
\textit{(But the market usually first notices those who can use tools and solve concrete tasks.)}

Kiến thức mạnh hơn rất nhiều khi nó đi kèm đôi tay biết làm.
\textit{(Knowledge becomes much stronger when it comes with hands that can do.)}

\end{truyen}

\begin{baihoc}
Giỏi thật là khi em nối được điều mình học với điều người ta đang cần.
\textit{(Real capability appears when you connect what you learned with what people actually need.)}
\end{baihoc}

\begin{ghinhoanh}
\textbf{Short English to remember:}

Real capability appears when you connect what you learned with what people actually need.

\end{ghinhoanh}

\begin{tuhocnhanh}
\textbf{Cau hoi tu hoc / Self-study questions}

1. Van de chinh la gi? \textit{What is the main problem?}

2. Bai hoc thuc te la gi? \textit{What is the practical lesson?}

3. Mot cau tieng Anh em muon nho la cau nao? \textit{Which English sentence do you want to remember?}
\end{tuhocnhanh}
